% 2010 年江苏卷第 17 题：电视塔与标杆测高示意。
% 来源 PDF：普通高考/2010/2010江苏.pdf，SHA-256 e701df878d51a700720549274fe7969bd54342f0f58acb051c516fc502b58af7。
% 原页：第 1 页（A3 三栏）第 3 栏上部，书页码 605。
% 重绘依据：A 为塔底，AE 为塔高 H；B 为标杆底，BC=h；D 在 B 左侧；
%   ∠ABE=α，∠ADE=β；DB 与 AB=d 标注。
\begin{tikzpicture}[semithick, scale=1.0, every node/.style={inner sep=1.6pt, font=\small}]
  \coordinate (D) at (0,0);
  \coordinate (B) at (2.2,0);
  \coordinate (A) at (5.2,0);
  \coordinate (C) at (2.2,0.8);
  \coordinate (E) at (5.2,5);

  \draw (D) -- (A);
  \draw (A) -- (E);
  \draw (B) -- (C);
  \draw[dashed] (D) -- (E);
  \draw[dashed] (B) -- (E);

  \node[below] at (D) {$D$};
  \node[below] at (B) {$B$};
  \node[below] at (A) {$A$};
  \node[left] at (C) {$C$};
  \node[right] at (E) {$E$};
  \node[below] at (3.7,0) {$d$};
  \node at (1.3,0.35) {$\beta$};
  \node at (3.0,0.35) {$\alpha$};
\end{tikzpicture}
